\subsection{Presentable rings}

A ring A is said to be presentable if there exists a regular ring R and a surjective homomorphism from R to A.

By definition, regular rings are presentable.

PROPOSITION 6.--- a) Every ring of fractions of a presentable ring is presentable. Every finitely generated algebra over a presentable ring is a presentable ring.

\begin{enumerate}
\def\labelenumi{\alph{enumi})}
\setcounter{enumi}{1}
\item
  Let A be a presentable ring and J an ideal of A. The separated completion \(\widehat{\mathbf{A}}\) of A for the J-adic topology is presentable.
\item
  Every complete Noetherian local ring is presentable.
\item
  Let A be a presentable local ring. There exist a regular local ring R and a surjective local homomorphism from R to A.
\end{enumerate}

Let A be a presentable ring; choose a regular ring R and a surjective homomorphism \(\rho: R \to A\) .

\begin{enumerate}
\def\labelenumi{\alph{enumi})}
\tightlist
\item
  Let S be a multiplicative subset of A; set \(T = \rho^{-1}(S)\) The homomorphism \(T^{-1}R \to S^{-1}A\) deduced from \(\rho\) is surjective and the ring \(T^{-1}R\) is regular (no. 3, example 1), hence \(S^{-1}A\) is presentable.
\end{enumerate}

Let B be a finitely generated A-algebra; there exists a finite set I and a surjective homomorphism \(\mathrm{A}[(\mathrm{T}_{i})_{i\in\mathbb{I}}]\to\mathrm{B}\) , hence a surjective homomorphism \(\mathrm{R}[(\mathrm{T}_{i})_{i\in\mathbb{I}}]\to\mathrm{B}\) . Since the ring \(\mathrm{R}[(\mathrm{T}_{i})_{i\in\mathbb{I}}]\) is regular (no. 2, cor. 5 of prop. 4), the ring B is presentable.

\begin{enumerate}
\def\labelenumi{\alph{enumi})}
\setcounter{enumi}{1}
\item
  Set \(1 = \rho^{-1}(\mathrm{J})\) and let \(\widehat{\mathbb{R}}\) denote the separated completion of \(\mathbb{R}\) for the I-adic topology; for each integer \(n \geqslant 0\) , the canonical map \(\mathrm{I}^n / \mathrm{I}^{n+1} \to \mathrm{J}^n / \mathrm{J}^{n+1}\) is surjective. Consequently, the homomorphism \(\widehat{\mathbb{R}} \to \widehat{\mathbb{A}}\) deduced from \(\rho\) is surjective (III, § 2, no. 8, cor. 2 of th. 1); since \(\widehat{\mathbb{R}}\) is regular (no. 2, cor. 3 of prop. 4), the ring \(\widehat{\mathbb{A}}\) is presentable.
\item
  This follows from IX, § 2, no. 5, th. 3 a) and IX, § 3, no. 3, th. 2 a).
\item
  Let \(\mathfrak{m}_{\Lambda}\) the maximal ideal of A; then \(\mathfrak{p} = \mathfrak{\rho}^{-1}(\mathfrak{m}_{\Lambda})\) is a prime ideal of R, the local ring \(R_{p}\) is regular and the homomorphism \(R_{p} \to A\) deduced from \(\mathfrak{p}\) is local and surjective.
\end{enumerate}

Since fields and Dedekind rings are regular, hence presentable, Proposition 6 implies that most rings commonly encountered in algebraic geometry are presentable.

PROPOSITION 7. - Let A be a presentable ring.

\begin{enumerate}
\def\labelenumi{\alph{enumi})}
\item
  The ring Λ is Noetherian and catenary.
\item
  Let M be a finitely generated A-module. The map

\[
\mathfrak {p} \mapsto \dim_ {\mathrm{A} _ {\mathfrak {p}}} (\mathrm{M} _ {\mathfrak {p}}) - \operatorname{prof} _ {\mathrm{A} _ {\mathfrak {p}}} (\mathrm{M} _ {\mathfrak {p}})
\]

from Spec(A) to Z is upper semicontinuous.

\item
  Let M be a \(\Lambda\) finitely generated module. The set of \(\mathfrak{p} \in \operatorname{Spec}(\Lambda)\) such that the \(\mathrm{A}_{\mathfrak{p}}\) -module \(\mathrm{M}_{\mathfrak{p}}\) being Macaulay is a dense open set. Its intersection with \(\operatorname{Supp}(\mathrm{M})\) is dense in \(\operatorname{Supp}(\mathrm{M})\) .
\end{enumerate}

Let us choose a regular ring R and a surjective homomorphism R → A.

\begin{enumerate}
\def\labelenumi{\alph{enumi})}
\item
  The ring R is Macaulay (no. 2, cor. 2 of prop. 4), hence A is catenary (§ 2, no. 2, prop. 2 and VIII, § 1, no. 3, rem. 2).
\item
  The R-module M is finitely generated and of finite projective dimension. \(< +\infty\) (no. 1, prop. 1). Let us identify \(\operatorname{Spec}(A)\) to a closed subset of \(\operatorname{Spec}(R)\) . Then the function \(\mathfrak{p} \mapsto \dim_{A_{\mathfrak{p}}} (M_{\mathfrak{p}}) - \operatorname{prof}_{A_{\mathfrak{p}}} (M_{\mathfrak{p}})\) on \(\operatorname{Spec}(A)\) is the restriction of the function \(\mathfrak{q} \mapsto \dim_{R_{\mathfrak{q}}} (M_{\mathfrak{q}}) - \operatorname{prof}_{R_{\mathfrak{q}}} (M_{\mathfrak{q}})\) on \(\operatorname{Spec}(R)\) It then suffices to apply Cor. 4 of Th. 1 of § 3, No. 5.
\item
  This is proved as in part c) of loc. cit.
\end{enumerate}

